\section{Zero-Knowledge Proof System}
\label{sec:zk}

\subsection{Circuit Description}

The withdrawal proof circuit $\mathcal{C}$ enforces four properties:

\begin{definition}[Privacy Pool ZK Circuit]
Public inputs: $(R, R_{\mathcal{A}}, \mathsf{nf}, \mathsf{recipient}, v)$\\
Private inputs: $(\mathsf{sk}, r, \mathsf{leafIndex}, \mathsf{merklePath}, \mathsf{assocPath})$

The circuit verifies:
\begin{enumerate}
    \item \textbf{Deposit existence}: $\mathsf{leaf} = H(\mathsf{Commit}(v, r), \mathsf{pk})$ is in the Merkle tree with root $R$ at path $\mathsf{merklePath}$.
    \item \textbf{Key ownership}: $\mathsf{pk} = \mathsf{sk} \cdot G$ (the prover knows the secret key).
    \item \textbf{Nullifier correctness}: $\mathsf{nf} = H_{\text{null}}(\mathsf{sk}, \mathsf{leafIndex})$.
    \item \textbf{Association set membership}: $\mathsf{leaf}$ is in the association set Merkle tree with root $R_{\mathcal{A}}$ at path $\mathsf{assocPath}$.
\end{enumerate}
\end{definition}

\subsection{Groth16 Proof System}

We use Groth16 \cite{groth2016size} over BLS12-381 for succinct proofs:

\begin{table}[H]
\centering
\caption{Groth16 proof parameters for the privacy pool circuit.}
\label{tab:groth16}
\begin{tabular}{lr}
\toprule
\textbf{Parameter} & \textbf{Value} \\
\midrule
Constraint count & $\sim$14,400 \\
Proof size & 192 bytes (3 $\mathbb{G}_1$ + 1 $\mathbb{G}_2$ elements) \\
Verification key size & $\sim$1.5 KB \\
Proving time (consumer HW) & 2.8 seconds \\
Verification time (on-chain) & 1.7 ms (precompile) \\
Verification gas & 326,000 \\
\bottomrule
\end{tabular}
\end{table}

\subsection{Circuit Constraints Breakdown}

\begin{table}[H]
\centering
\caption{Constraint breakdown by circuit component.}
\label{tab:constraints}
\begin{tabular}{lrr}
\toprule
\textbf{Component} & \textbf{Constraints} & \textbf{Share} \\
\midrule
Commitment hash (Poseidon) & 312 & 2.2\% \\
Nullifier derivation (Poseidon) & 624 & 4.3\% \\
Merkle path (depth 20) & 6,552 & 45.7\% \\
Association set path (depth 20) & 6,864 & 47.8\% \\
\midrule
\textbf{Total} & \textbf{14,352} & \textbf{100\%} \\
\bottomrule
\end{tabular}
\end{table}

\subsection{Hash Function Selection}

We use Poseidon \cite{grassi2021poseidon} as the in-circuit hash function, optimized for arithmetic circuits:

\begin{definition}[Poseidon Hash]
Poseidon is an algebraic hash function over $\mathbb{F}_p$ (the BLS12-381 scalar field) with:
\begin{itemize}
    \item Width: $t = 3$ (2 inputs + 1 capacity)
    \item Full rounds: $R_f = 8$
    \item Partial rounds: $R_p = 57$
    \item Security: 128-bit against algebraic and statistical attacks
\end{itemize}
Poseidon requires $\sim$300 R1CS constraints per hash, versus $\sim$25,000 for SHA-256 in a circuit.
\end{definition}

